% BACKMATTER:VERSION:START
\section*{Version Statement}

This is version~0.10.5 of the manuscript, dated
2026-11-26, prepared for Zenodo
deposit as a preprint.
% BACKMATTER:VERSION:END

\section*{Data availability}
\addcontentsline{toc}{section}{Data availability}

This document is a normative specification; no empirical data about deployments is
used; the measurements in Section~\ref{sec:conformance} are synthetic. The evidence is executable and is deposited with the paper in a
single Zenodo record: the PDF and an archive of the source tree at the
release tagged \texttt{v0.10.5}. A copy of the paper's source without these
artefacts may be posted elsewhere (for example on arXiv); the artefacts are
then reached through the Zenodo record and the repository named below, not
through that copy. The archive contains the conformance suite and Hypothesis
properties (T1--T158), the mutation check, the test vectors and their runner,
the Rust crate, the TLA+ model with its TLC configurations, and the event
schemas. The same tree is published at
\url{https://github.com/strebulaev-articles/articles.bridge_contract_for_agent_session_governance.public};
the Zenodo record is the citable snapshot of tag \texttt{v0.10.5}. The paper text is released under
CC-BY-4.0 and the code under the MIT licence (files \texttt{LICENSE} and
\texttt{LICENSE-CODE}).

The \texttt{paper/sources/} directory of the archive (it is not part of a copy
of the paper posted without the artefacts) holds the source
texts that were read for the bibliography and the mapping tables, with a
manifest (\texttt{paper/sources/MANIFEST.md}) listing the URL and the
SHA-256 of each version cited. Those source texts are the IETF
Internet-Drafts, the arXiv preprint, and the Herdr documentation and
licence; they are reproduced here only as the working copies that were
verified by hash, and they remain the property of their respective
authors under their respective licences (the IETF Trust's provisions for
Internet-Drafts, the arXiv non-exclusive licence, the Apache-2.0
licence for Herdr). They are not deposited under the CC-BY-4.0 or MIT
terms that cover the rest of this archive.

To reproduce the reported results from the archive, run in
\texttt{paper/supplementary}: \texttt{python3 -I -m pytest
conformance\_suite.py -q}, \texttt{python3 -I mutation\_check.py},
\texttt{python3 -I run\_vectors.py -{}-check-paper
../sections/generated\_vectors.tex}, and \texttt{cargo test} in
\texttt{rust-bridge}. The model-checking commands and configurations
are in \texttt{formal-spec/README.md}.

\section*{AI use}
\addcontentsline{toc}{section}{AI use}

AI tools were used during the preparation of this manuscript for
grammatical correction and technical formatting (sentence structure for
academic tone, formatting of the bibliographic references), for code
generation, and for review of the text and the citations. The
supplementary code (conformance suite, mutation check, vector runner,
Rust crate, TLA+ model and event schemas) was written with AI
assistance. The bibliography was checked against the cited sources by
the human author, who read each cited draft at the version recorded in
\texttt{paper/sources/MANIFEST.md} and verified the cited section
numbers against the text. One entry (the companion
\cite{I-D.hood-agtp-standard-methods} vocabulary, which expired on
2026-10-18) is cited on the strength of an automated retrieval that
confirmed the draft's title and section structure but did not yield a
hashable byte stream; this limitation is recorded in the manifest, and
rows in the comparison table (Table~\ref{tab:comparison}) that depend
on that entry are marked accordingly. AI was not employed for
literature synthesis, data analysis, or the generation of the core
ideas, hypotheses, and scientific conclusions. The human author
retains full responsibility for the content; they reviewed the edited
text to ensure terminological precision and performed bibliographic
verification of the entries the manifest records.

\section*{Competing interests}
\addcontentsline{toc}{section}{Competing interests}

The author develops Termenbox, the proprietary, free-to-use, self-hosted web terminal for
AI coding agents to which the contract was mapped (Section~\ref{sec:impl}
and the mapping in Section~\ref{sec:mapping}) and which may adopt the
specification: adoption is intended and not scheduled. The specification is written to be implementation-agnostic;
the interest is disclosed so that readers can weigh the Termenbox findings
accordingly.
